\documentclass[a4paper,12pt]{article}
\usepackage[T2A]{fontenc}
\usepackage[utf8]{inputenc}
\usepackage[russian]{babel}
\usepackage{amsmath,amssymb,mathtools}
\renewcommand{\sfdefault}{cmss}
\renewcommand{\familydefault}{\sfdefault}
\usepackage{icomma}
\usepackage[left=23mm,right=23mm,top=23mm,bottom=22mm]{geometry}
\usepackage{microtype}
\usepackage{tikz}
\usetikzlibrary{arrows.meta,positioning,shapes.geometric,calc}
\usepackage{booktabs,tabularx}
\usepackage{enumitem}
\usepackage{fancyhdr}
\usepackage{setspace}
\usepackage{xcolor}
\usepackage[hidelinks]{hyperref}

\definecolor{accent}{HTML}{126E71}
\definecolor{darkaccent}{HTML}{26494C}
\definecolor{textgray}{HTML}{455056}
\definecolor{lightline}{HTML}{CAD9D7}
\setstretch{1.28}
\setlength{\parindent}{0pt}
\setlength{\parskip}{3pt}
\emergencystretch=2.5em
\setlist[enumerate]{leftmargin=1.8em,itemsep=2pt,topsep=2pt}
\setlist[itemize]{leftmargin=1.5em,itemsep=1pt,topsep=1pt}
\pagestyle{fancy}
\fancyhf{}
\fancyhead[L]{\small\color{textgray}АННА БАННОВА\quad /\quad 18.10.26}
\fancyhead[R]{\small\color{textgray}МАТЕМАТИКА\quad /\quad ЕГЭ}
\fancyfoot[C]{\small\color{textgray}\thepage}
\renewcommand{\headrulewidth}{0.35pt}
\renewcommand{\footrulewidth}{0pt}
\setlength{\headheight}{15pt}
\newcommand{\sect}[1]{\vspace{3pt}{\large\bfseries\color{darkaccent}#1}\par\vspace{2pt}\noindent\color{lightline}\rule{\linewidth}{0.5pt}\color{black}\vspace{2pt}}
\newcommand{\keyphrase}[1]{{\color{accent}\bfseries #1}}
\newcommand{\num}[1]{\text{#1}}

\begin{document}
\thispagestyle{empty}
\begin{center}
\vspace*{35mm}
{\small\bfseries\color{accent} МАТЕМАТИКА \quad / \quad ПРОФИЛЬНЫЙ ЕГЭ}\par
\vspace{14mm}
{\fontsize{31}{38}\selectfont\bfseries\color{darkaccent}Чек-лист}\par
\vspace{6mm}
{\fontsize{19}{26}\selectfont Прикладные задачи со степенями\\и большими числами}\par
\vspace{15mm}
{\color{lightline}\rule{100mm}{0.7pt}}\par
\vspace{13mm}
{\large Анна Баннова}\par
\vspace{3mm}
{\normalsize 18 октября 2026 года}\par
\vspace{20mm}
{\small\color{textgray}Стандартная запись числа, разложение на множители,\\степени с дробными показателями, преобразование формул}\par
\vfill
{\small\bfseries\color{darkaccent}Лёвин Артём Александрович эксклюзивно для Анны Банновой}\par
\vspace{8mm}
\begin{tikzpicture}[x=1cm,y=1cm]
\draw[accent!55,line width=1.1pt] (-7,0) .. controls (-3.9,0.55) and (-1.4,-0.35) .. (0.5,0.15)
.. controls (2.7,0.60) and (4.5,-0.13) .. (7,0.22);
\end{tikzpicture}
\end{center}
\clearpage

\sect{01\quad Идея занятия}
В прикладной задаче неизвестная величина часто находится в степени. Прежде чем извлекать корень, удобно \keyphrase{привести большие числа к произведениям степеней простых множителей} и сократить общие множители. Этот приём значительно упрощает расчёты.

\sect{02\quad Как обращаться с большими числами}
Изменение записи числа компенсируется изменением показателя десятки:
\[
1{,}5625\cdot10^{25}=15625\cdot10^{21},\qquad
6{,}25\cdot10^5=625\cdot10^3.
\]
При переносе запятой \emph{вправо} на $m$ раз показатель десятки уменьшается на $m$; при переносе \emph{влево} увеличивается на $m$.

\textbf{Разложение на простые множители:}
\[
\begin{aligned}
15625&=5^6,\\
228&=2^2\cdot3\cdot19,\\
57&=3\cdot19,\\
625&=5^4,\qquad 10=2\cdot5.
\end{aligned}
\]
Например, разложение позволяет сократить целую дробь, не вычисляя произведения больших чисел:
\[
\frac{228}{57}=\frac{2^2\cdot3\cdot19}{3\cdot19}=2^2=4.
\]

\sect{03\quad Свойства степеней}
При положительных основаниях $a,b$ и допустимых показателях $m,n$:
\[
\begin{aligned}
a^m\cdot a^n&=a^{m+n},&\quad \frac{a^m}{a^n}&=a^{m-n},\\
(a^m)^n&=a^{mn},& (ab)^n&=a^n b^n.
\end{aligned}
\]
\textbf{При делении степеней десятки вычитаем весь показатель знаменателя:}
\[
\frac{10^{23}}{10^{20}\cdot10^{-9}}
=10^{23-20-(-9)}=10^{12}.
\]
\textbf{Как убрать степень у неизвестной:} для $x>0$ и $r\ne0$ из $x^r=A>0$ следует $x=A^{1/r}$.
\[
T^4=A\ (T>0)\quad\Longrightarrow\quad T=A^{1/4};
\qquad V^{5/3}=B\ (V>0)\quad\Longrightarrow\quad V=B^{3/5}.
\]

\clearpage
\sect{04\quad Модель I. Излучение и температура}
\textbf{Формула:} $P=\sigma ST^4$, где $P$~--- мощность излучения, $\sigma$~--- коэффициент, $S$~--- площадь поверхности, $T>0$~--- абсолютная температура.

\textbf{Данные разобранного примера:}
\[
P=1{,}5625\cdot10^{25},\qquad
\sigma=5{,}7\cdot10^{-8},\qquad
S=\frac{10^{20}}{228}.
\]
Требуется найти температуру, соответствующую данному значению мощности.

\textbf{Шаг 1. Приводим числа к удобной форме и подставляем:}
\[
P=15625\cdot10^{21},\qquad \sigma=57\cdot10^{-9},
\]
\[
15625\cdot10^{21}=57\cdot10^{-9}\cdot\frac{10^{20}}{228}\cdot T^4.
\]

\textbf{Шаг 2. Выражаем четвёртую степень температуры:}
\[
T^4=\frac{15625\cdot10^{21}\cdot228}{57\cdot10^{-9}\cdot10^{20}}.
\]

\textbf{Шаг 3. Раскладываем числа, сокращаем и объединяем степени:}
\[
\begin{aligned}
T^4
&=\frac{5^6\cdot10^{21}\cdot 2^2\cdot3\cdot19}
{3\cdot19\cdot10^{-9}\cdot10^{20}}\\[3pt]
&=5^6\cdot2^2\cdot10^{21-(-9)-20}\\
&=5^6\cdot2^2\cdot10^{10}\\
&=5^6\cdot2^2\cdot(2\cdot5)^{10}\\
&=5^{16}\cdot2^{12}.
\end{aligned}
\]

\textbf{Шаг 4. Возводим обе части в степень $\frac14$:}
\[
T=(5^{16}\cdot2^{12})^{1/4}=5^4\cdot2^3=625\cdot8
=\boxed{5000}.
\]

\keyphrase{Контроль:} $T=5000>0$; подстановка в $P=\sigma ST^4$ возвращает заданную мощность.

\textbf{Если условие содержит слова «мощность не менее $P_0$»:} сначала находят \emph{граничную температуру} из равенства $P=P_0$. Так как $P=\sigma ST^4$ возрастает при $T>0$, требованию $P\ge P_0$ отвечает $T\ge T_0$.

\clearpage
\sect{05\quad Модель II. Давление и объём}
\textbf{Формула:} $pV^k=C$, где $p>0$~--- давление, $V>0$~--- объём, $k>0$~--- заданный показатель, $C>0$~--- постоянная.

\textbf{Данные разобранного примера:}
\[
k=\frac53,\qquad C=200,\qquad p=6{,}25\cdot10^5.
\]
Найдём объём при данном значении давления.

\textbf{Шаг 1. Переписываем давление и подставляем:}
\[
p=625\cdot10^3,\qquad
(625\cdot10^3)V^{5/3}=200.
\]

\textbf{Шаг 2. Выражаем степень объёма и сокращаем:}
\[
\begin{aligned}
V^{5/3}
&=\frac{200}{625\cdot1000}
=\frac{1}{625\cdot5}\\
&=\frac1{5^4\cdot5}=\frac1{5^5}=5^{-5}.
\end{aligned}
\]

\textbf{Шаг 3. Возводим обе части в обратную степень $\frac35$:}
\[
\begin{aligned}
V&=(5^{-5})^{3/5}=5^{-5\cdot3/5}
=5^{-3}\\
&=\frac1{125}=\boxed{0{,}008}.
\end{aligned}
\]
\keyphrase{Проверка:} $V^{5/3}=(5^{-3})^{5/3}=5^{-5}$, а
\[
(625\cdot1000)\cdot5^{-5}=(5^4\cdot10^3)\cdot\frac1{5^5}=200.
\]

\textbf{Если требуется выполнить условие $p\ge p_0$:} при постоянных $C>0$ и $k>0$ давление $p=C/V^k$ \emph{убывает} с ростом объёма. Поэтому после нахождения границы $V_0$ из $p=p_0$ получаем $0<V\le V_0$.

\sect{06\quad Типичные ошибки и самопроверка}
\textbf{Показатели и сокращение.} Запишите $23-20-(-9)=12$, не теряя минус перед отрицательным показателем. Помните, что $57=3\cdot19$, а $(2\cdot5)^{10}=2^{10}\cdot5^{10}$.

\textbf{Обратная степень.} Для $V^{5/3}=A$ используйте $V=A^{3/5}$, поскольку $\frac53\cdot\frac35=1$. В конце проверьте вычисление $\frac1{125}=0{,}008$ обратным умножением.

\clearpage
\sect{07\quad Алгоритм: прикладная задача со степенями}
\vspace{8mm}
\begin{center}
\begin{tikzpicture}[
 >={Stealth[length=2.7mm]}, line width=.8pt, font=\small,
 start/.style={ellipse,draw=accent,minimum height=10mm,text width=67mm,align=center,inner sep=4pt},
 action/.style={rectangle,rounded corners=2pt,draw=lightline!70!black,
 text width=95mm,align=center,minimum height=12mm,inner sep=5pt},
 arr/.style={->,draw=darkaccent}]
\node[start] (a) at (0,0) {Прочитать формулу и условие};
\node[action] (b) at (0,-2.05) {Выписать данные, искомую величину и возможное ограничение};
\node[action] (c) at (0,-4.1) {Перевести числа в удобную запись: степенная форма, множители};
\node[action] (d) at (0,-6.15) {Подставить значения и выразить степень неизвестной};
\node[action] (e) at (0,-8.2) {Разложить составные числа, сократить общие множители};
\node[action] (f) at (0,-10.25) {Собрать одинаковые основания, упростить степени};
\node[action] (g) at (0,-12.3) {Возвести обе части в обратную степень; вычислить результат};
\node[action] (h) at (0,-14.35) {Проверить подстановку, физический смысл и направление ограничения};
\node[start] (i) at (0,-16.4) {Записать ответ с нужными единицами};
\foreach \x/\y in {a/b,b/c,c/d,d/e,e/f,f/g,g/h,h/i} {\draw[arr] (\x.south) -- (\y.north);}
\end{tikzpicture}
\end{center}
\end{document}
